%% kapa: exam=1 round=21 year=2010 period=1 form=A subject=economics subject_ko=경제원론 seq=31 number=71
%% source: Q-Net 감정평가사 기출문제 (cst00309) artlSeq=5024181 file=extracted/1008/5024181/2010년 제21회 감정평가사 1차시험 기출문제/1교시.hwp
%% conversion: kordoc 4.24.51 -> tools/gen1.py (auto); review=proofread-sonnet; flags=none
\begin{question}{71}
A기업의 생산함수는 Y = L - 100이고, 노동 1단위당 임금은 1이다. 다음 설명 중 옳은 것은? (단, L은 노동, Y는 생산량, Y > 0이다.)
\begin{choices}
  \item 노동의 한계생산이 체감한다.
  \item 노동의 평균생산은 일정하다.
  \item 생산량이 늘어남에 따라 평균비용은 처음에는 감소하나 생산량이 일정 수준을 넘어서면 점차 증가한다.
  \item 생산량이 일정 수준을 넘어서면 한계비용이 평균비용 보다 더 커진다.
  \item 규모의 경제가 나타난다.
\end{choices}
\end{question}
\answer{5}
